
This work sits at the intersection of three research areas: LLM self-repair with feedback oracles, specification grounding for code generation, and statistical design for paired LLM repair evaluation.

\subsubsection{2.1 LLM Self-Repair and Feedback Oracles}

Self-Refine \citep{Madaan2023SelfRefine} established that iterative self-generated feedback improves performance by approximately 20 percentage points across multiple tasks. The Self-Refine setting conflates re-exposure with informative self-generated feedback, making it difficult to isolate the contribution of feedback content from the contribution of additional sampling. The blind reprompting condition in the present design (Condition B) reduces Self-Refine's iterative self-feedback to a single fixed reprompt string, isolating re-exposure from informative content.

FeedbackEval \citep{Dai2025FeedbackEval} provides the closest systematic study of feedback type effects on code repair, finding that mixed feedback achieves a 63.6\% fix rate compared to 53.1\% for minimal feedback on a purpose-built benchmark. Critically, removing docstring context from feedback prompts causes severe performance degradation, providing direct empirical motivation for the docstring component of the specification triple. However, FeedbackEval does not isolate specification context from blind re-exposure through a paired design, nor does it evaluate on EvalPlus.

Arimbur (2026) evaluates iterative self-repair on HumanEval across model scales, finding improvement of 4.9 to 17.1 percentage points over several rounds, with assertion errors proving hardest at approximately 45\% failure rate. This confirms that EvalPlus-style failures are difficult for self-repair but does not test structured specification context as an oracle.

The DUALFIX pipeline \citep{Akli2026Failing} demonstrates that execution-feedback repair leaves a residual failure class addressable only by specification-level intervention, with evolved natural language repair rules fixing 12--17\% of cases that execution repair cannot address. This finding directly supports the claim that semantic-level context is necessary beyond error traces.

\subsubsection{2.2 Specification Grounding for Code Generation}

Haeri and Ghelichi (2026) show that specification grounding improves correct code generation by 38 percentage points, with docstring context identified as the primary driver. Although the result is on code generation rather than code repair, the underlying mechanism --- docstring re-anchoring preventing goal misspecification --- transfers directly to the repair setting, where the model may regenerate the same incorrect solution without explicit re-anchoring on the intended algorithm.

ContrastRepair \citep{Kong2024ContrastRepair} tests contrastive input/output pairs for automated program repair on the Defects4J corpus, finding 143 of 337 bugs fixed with contrastive pairs compared to 124 with single failure messages. The contrastive pair design --- providing expected output alongside actual output --- is directly analogous to the input/output pair component of the specification triple. The key difference is that ContrastRepair targets Java bugs in a classical automated program repair setting, not LLM code generation on EvalPlus, and does not include docstring re-anchoring as a third component.

VRpilot \citep{Kulsum2024Case} applies chain-of-thought reasoning with patch validation to vulnerability repair, finding 14\% more correct patches than naive repair. The consistent advantage of structured semantic context over minimal feedback across repair domains provides convergent support for the present design.

SGCR \citep{Wang2025SGCR} validates specification-grounded code review at industry scale, with a 42\% developer adoption rate and 90.9\% improvement in review utility, establishing that specification grounding has practical value in production code assessment settings.

\subsubsection{2.3 Statistical Design for LLM Repair Evaluation}

Iscan (2026) provides the methodological foundation for the McNemar design used here. The paper validates placebo-controlled decomposition of self-repair feedback using frozen small code models, demonstrating that content (code plus facts) outperforms bare code re-exposure by 18 percentage points (p = 0.00042). The paper explicitly validates the one-tailed McNemar design for paired LLM repair comparisons --- the same design pre-registered in the present work for predictions P1 and P2.

The EvalPlus benchmark \citep{Liu2023EvalPlus} is the evaluation environment for this work. Its $80\times$ test augmentation over base HumanEval produces a substantially stricter evaluation that specifically surfaces semantic divergences surviving round-0 evaluation --- the failure population of greatest interest for specification-aligned repair.

\subsubsection{2.4 Position Relative to Prior Work}

Relative to the above literature, the present work contributes (1) an explicit four-condition verification of the EvalPlus failure set as a reproducible data source for repair experiments, (2) the first three-condition McNemar design isolating specification-aligned repair from blind reprompting on EvalPlus's own semantic failure set under its augmented test suite, and (3) pre-registration of comparative predictions before mechanism experiment execution.

